\scanpage{80}
% lines 1–5, qualification at upper edge and unit condition
\inserted{scindage $x_0$ d'un $P_0$ \unsure{muni déjà|muni} d'une} $\unsure{\chi|\gamma}$-trivialisation est ``\uline{unitaire}'', si l'élém
\[
b_0=\pi(x_0)
\]
de $L$ est \deleted{\unsure{déf}} \unsure{tel|telle} que $\{b\}$ soit une \uline{base} de $L$.
\deleted{Donc} Donc grâce à cette base, $L\xleftarrow{\sim}\underline O_X,\ b=1$, et on aura donc
% lines 6–10, coordinates for the unit
\[
\begin{aligned}
P_0&\simeq\underline O_{\deleted{X}},\quad x_0=1\\
E_0&=\underline O\times\underline O=\{(u,\lambda)\mid u,\lambda\in\Gamma O\}\xrightarrow{\psi}\lambda
\end{aligned}
\]
et
\[
\left|\begin{aligned}
\pi(u,\lambda)&=\lambda+\chi u\\
T&=(-1,\chi)
\end{aligned}\right.
\]
% lines 11–15, unit isomorphisms
Ceci posé, si $\deleted{E}P_1$ est un torseur $\chi$-trivialisé, alors les
\[
\left.\begin{aligned}
P_0*P_1&\simeq P_1\\
\underset{\shortparallel\,\underline O}{L_0}\otimes L_1&\simeq L_1\\
x_1&\longmapsto x_0*x_1\\
u_1&\mathbin{\unsure{\mapsto|-}}b_0\otimes\deleted{\unsure{x}}u_1
\end{aligned}\right.
\]
% lines 16–18, compatibility sentence
Les compatibilités entre les contraintes d'unité, d'assoc., de commutativité... .
% lines margin beside lines 11–20, diagonal writing from lower left
% slanted
\margin{\uline{NB} Un $(E,\psi,T)$ est \uline{inversible} pour $*$ sss\\
a) $L$ inversible\\
b) Sur \inserted{$X_0={}$}$V(\chi)\subset X$ (le schéma des zéros de $\chi$), la section $T_0=T|X_0$ de $({}_{\chi}P)_0\simeq{}_{\chi_0}P_0\simeq L_0$ est inversible. \deleted{\unsure{de}}\\
Si $(E,\psi,T)$ est scindé, et si $b\in\Gamma L$ correspond à un scindage $x_0$, \unsure{cf. devant|cf. dessus|alors} (on a $T=2\chi x-b$) donc $T_0=-b_0$, et la condition est que $b_0$ soit section inv de $L_0$ i.e. \unsure{qu’il en} soit une base. Cela équivaut à la condition $\unsure{L^{\otimes2}=2L^{\otimes2}+b\cdot L}$ \unsure{, ou à} $L',\unsure{L^{-\otimes}|L^{-1}}$ \deleted{\unsure{si}}\\
La condition d'inversibilité équivaut aussi à celle-ci\\
a) $L$ inversible\\
b) $\exists$ \uline{localement} un scindage $x$ de $(E,\psi,T)$, tel que le $b$ correspondant soit une base de $L$.}
% lines 19–21, new heading and specialization
\centerline{---}
Relation avec les involutions canoniques.

On suppose à partir de maintenant
\[
\chi=2.
\]
% archive sheet number: 79
